\section*{अध्याय \ref{ch:b1:precipitation} --- अवक्षेपण और विलेयता}

\begin{solution}{exo:b1:precipitation:1}
\ce{BaSO4(s) <=> Ba^2+ + SO4^2-} के लिए $K_s = [\ce{Ba^2+}][\ce{SO4^2-}] =
10^{-9.97}$। \ce{CaF2(s) <=> Ca^2+ + 2F-} के लिए $K_s = [\ce{Ca^2+}][\ce{F-}]^2 =
10^{-9.84}$। \ce{Ag2CrO4(s) <=> 2Ag+ + CrO4^2-} के लिए $K_s =
[\ce{Ag+}]^2[\ce{CrO4^2-}] = 10^{-11.95}$। \ce{Fe(OH)3(s) <=> Fe^3+ + 3OH-} के लिए
$K_s = [\ce{Fe^3+}][\ce{OH-}]^3 = 10^{-38.55}$। $K_s = \num{2.0e-13}$ के लिए
$\mathrm{p}K_s = 12.70$।
\end{solution}

\begin{solution}{exo:b1:precipitation:2}
$s = \sqrt{K_s} = 10^{-12.27/2} = \qty{7.3e-7}{mol/L}$, अर्थात्
$\num{7.3e-7} \times 187.8 = \qty{1.4e-4}{g/L} = \qty{0.14}{mg/L}$।
\end{solution}

\begin{solution}{exo:b1:precipitation:3}
मिलाने के बाद (आयतन दोगुना), $[\ce{Ba^2+}] = \qty{1.0e-3}{mol/L}$ और
$[\ce{SO4^2-}] = \qty{1.0e-4}{mol/L}$: $Q_s = \num{1.0e-7} > K_s =
\num{1.1e-10}$, बेरियम सल्फ़ेट अवक्षेपित होता है। दूसरी स्थिति में
$[\ce{SO4^2-}] = \num{1.0e-7}$ और $Q_s = \num{1.0e-10} < K_s$: कोई
अवक्षेप नहीं, यद्यपि बाल-बाल।
\end{solution}

\begin{solution}{exo:b1:precipitation:4}
पानी में $s = 10^{-9.97/2} = \qty{1.0e-5}{mol/L}$। \qty{0.010}{mol/L}
सल्फ़ेट में (\cref{prop:b1:precipitation:common-ion}), $s = 10^{-9.97}/0.010 =
\qty{1.1e-8}{mol/L}$, लगभग एक हज़ार गुना कम (गुणक 970)।
\end{solution}

\begin{solution}{exo:b1:precipitation:5}
पानी में $s = (10^{-9.84}/4)^{1/3} = \qty{3.3e-4}{mol/L}$, अर्थात्
\qty{26}{mg/L}। \qty{0.10}{mol/L} कैल्सियम क्लोराइड में
$[\ce{Ca^2+}] \approx 0.10$ और $[\ce{F-}] = 2s$, इसलिए
$s = \frac12\sqrt{10^{-9.84}/0.10} = \qty{1.9e-5}{mol/L}$।
सन्निकटन $s \ll 0.10$ परिमाण की चार कोटियों से सही है।
\end{solution}

\begin{solution}{exo:b1:precipitation:6}
आरंभ: $\mathrm{pH} = 14.00 - \frac12(11.25 + \log 0.050) = 14.00 - 4.97 =
9.03$। \qty{99}{\percent} पर $[\ce{Mg^2+}] = \num{5.0e-4}$ और
$\mathrm{pH} = 14.00 - \frac12(11.25 - 3.30) = 10.03$। pH अक्ष पर
\ce{Mg(OH)2}, 9.03 से ऊपर उपस्थित है; उसके नीचे मैग्नीशियम पूरी तरह \ce{Mg^2+} है,
\qty{0.050}{mol/L} पर।
\end{solution}

\begin{solution}{exo:b1:precipitation:7}
सिल्वर ब्रोमाइड पहले प्रकट होता है, $[\ce{Ag+}] = 10^{-12.27 + 2} =
10^{-10.27}$ पर, सिल्वर क्लोराइड $10^{-7.75}$ पर। जब सिल्वर क्लोराइड
प्रकट होता है, तब $[\ce{Br-}] = 10^{-12.27}/10^{-7.75} = 10^{-4.52} =
\qty{3.0e-5}{mol/L}$, ब्रोमाइड का \qty{0.30}{\percent}। \qty{0.1}{\percent} की सामान्य
कसौटी पर पृथक्करण स्वीकार्य नहीं है: दोनों
$\mathrm{p}K_s$ में केवल 2.5 का अंतर है।
\end{solution}

\begin{solution}{exo:b1:precipitation:8}
सिल्वर क्रोमेट तब बनता है जब $[\ce{Ag+}]^2 \times \num{5.0e-3} = 10^{-11.95}$,
$[\ce{Ag+}] = \qty{1.5e-5}{mol/L}$। तब $[\ce{Cl-}] = 10^{-9.75}/\num{1.5e-5}
= \qty{1.2e-5}{mol/L}$, क्लोराइड का \num{2.4e-4} (\qty{0.024}{\percent}) भिन्न:
लाल रंग तब प्रकट होता है जब क्लोराइड व्यावहारिक रूप से पूरा
अवक्षेपित हो चुका होता है।
\end{solution}

\begin{solution}{exo:b1:precipitation:9}
$[\ce{CO3^2-}] = [\ce{HCO3-}]K_{a2}/[\ce{H3O+}] = [\ce{HCO3-}]\,10^{\mathrm{pH}
- \mathrm{p}K_{a2}}$, इसलिए $K_s = [\ce{Ca^2+}][\ce{HCO3-}]\,10^{\mathrm{pH} -
\mathrm{p}K_{a2}}$, जो माँगा गया संबंध है। $[\ce{Ca^2+}] =
[\ce{HCO3-}] = s$ के साथ: $s^2 = 10^{-8.30 + 10.33 - 8.30} = 10^{-6.27}$,
$s = \qty{7.3e-4}{mol/L}$ (\ce{CaCO3} का \qty{73}{mg/L}), जबकि
यदि कार्बोनेट आयन क्षारक न होता तो $10^{-4.15} = \qty{7.1e-5}{mol/L}$।
\end{solution}

\begin{solution}{exo:b1:precipitation:10}
1.~\ce{CO2(g) <=> CO2(aq)} जोड़ने पर: $K' = K K_H = 10^{-4.34 - 1.47} =
10^{-5.81}$।
2.~$[\ce{Ca^2+}] = s$ और $[\ce{HCO3-}] = 2s$, इसलिए $4s^3 = K'p$।
\qty{0.010}{bar} पर: $s = (10^{-5.81} \times 0.010/4)^{1/3} =
\qty{1.6e-3}{mol/L}$, \qty{157}{mg/L}। बाहरी वायु में:
$s = \qty{5.5e-4}{mol/L}$, \qty{55}{mg/L}। (जाँच: मिट्टी की वायु में
$[\ce{CO2(aq)}] = 10^{-3.47}$ और $\mathrm{pH} = 6.37 + \log(3.14\times
10^{-3}/3.39 \times 10^{-4}) = 7.34$: हाइड्रोजनकार्बोनेट सचमुच प्रभावी है।)
3.~मिट्टी के ऊँचे कार्बन डाइऑक्साइड दाब में कैल्साइट से संतृप्त हुआ पानी
गुफा की छत तक पहुँचता है, जहाँ दाब कम है: वह
कार्बन डाइऑक्साइड खो देता है, $Q > K'$, और जहाँ बूँद लटकती है वहाँ कैल्साइट
अवक्षेपित होता है, वलय पर वलय।
% ledger: air:CO2, dfg:CO2_g, dfg:CO2_ao (log KH computed in figdata/bachelor-1/aqueous-constants.py)
\end{solution}

\begin{solution}{exo:b1:precipitation:11}
1.~आरंभ $\mathrm{pH} = 14.00 - \frac12(16.38 - 2.00) = 6.81$ पर। पूर्ण
पुनर्विलयन तब, जब $[\ce{Zn(OH)4^2-}] = 10^{-2} = 10^{-1.93}[\ce{OH-}]^2$,
$[\ce{OH-}] = 10^{-0.035}$, $\mathrm{pH} = 13.97$। हाइड्रॉक्साइड
pH 6.81 और 13.97 के बीच उपस्थित है।
2.~$\mathrm{pH} = 14.00 - 14.45/4 = 10.39$, $s_{\min} = 2 \times
10^{-(16.38 + 1.93)/2} = \qty{1.4e-9}{mol/L}$ (इस दो-स्पीशीज़ मॉडल में)।
3.~प्रवणता $-2$ वाली रेखा, $(6.81, -2)$ से नीचे न्यूनतम तक, जो
$(10.39, -8.9)$ के पास है, फिर प्रवणता $+2$ से ऊपर $(13.97, -2)$ तक; वास्तविक वक्र
अन्य हाइड्रॉक्सो स्पीशीज़ के कारण गोलाई लिए होता है
(\cref{sec:b1:precipitation:amphoteric} का चित्र)।
\end{solution}

\begin{solution}{exo:b1:precipitation:12}
1.~$\ce{Al(OH)3(s) + OH- <=> Al(OH)4-}$, $K = 10^{-1.23}$; सारा
ऐलुमिनियम तब घुला होता है जब $10^{-3} \leqslant K[\ce{OH-}]$, अर्थात्
$[\ce{OH-}] \geqslant 10^{-1.77}$, $\mathrm{pH} \geqslant 12.23$।
2.~$[\ce{Fe^3+}] = 10^{-38.55}/10^{-3 \times 1.77} = 10^{-33.2}$: आयरन
पूरी तरह \ce{Fe(OH)3} के रूप में रहता है; छानने से आयरन (ठोस)
ऐलुमिनियम (निस्यंद) से अलग हो जाता है।
3.~मैग्नीशियम हाइड्रॉक्साइड उभयधर्मी नहीं है: आधिक्य हाइड्रॉक्साइड में वह
आयरन(III) हाइड्रॉक्साइड के साथ ठोस ही रहता है, और कुछ अलग नहीं होता।
\end{solution}

\begin{solution}{pb:b1:precipitation:1}
\textbf{1.} \ce{Fe(OH)3(s) <=> Fe^3+ + 3OH-}, $K_s = [\ce{Fe^3+}][\ce{OH-}]^3$;
\ce{Zn(OH)2(s) <=> Zn^2+ + 2OH-}, $K_s = [\ce{Zn^2+}][\ce{OH-}]^2$;
\ce{Mg(OH)2(s) <=> Mg^2+ + 2OH-}, इसी प्रकार।
\textbf{2.} स्थिरांकों का रूप एक-सा नहीं है ($[\ce{OH-}]$ की घात 3 या 2)
और तीनों आयन भिन्न सांद्रताओं पर हैं; केवल
आरंभ के pH की तुलना की जा सकती है।
\textbf{3.} आरंभ पर $c[\ce{OH-}]^n = K_s$, इसलिए $n\log[\ce{OH-}] =
-\mathrm{p}K_s - \log c$ और $\mathrm{pH} = \mathrm{p}K_e +
\log[\ce{OH-}] = \mathrm{p}K_e - \frac1n(\mathrm{p}K_s + \log c)$।
\textbf{4.} $14.00 - \frac13(38.55 - 3.00) = 2.15$।
\textbf{5.} $14.00 - \frac12(16.38 - 2.30) = 6.96$।
\textbf{6.} $14.00 - \frac12(11.25 - 1.70) = 9.22$।
\textbf{7.} आयरन(III) (2.15), फिर ज़िंक (6.96), फिर मैग्नीशियम (9.22)।
pH 2.0 पर आयरन के लिए $Q_s = 10^{-3} \times 10^{-36} = 10^{-39} < 10^{-38.55}$,
और बाक़ी अवक्षेपण से और भी दूर हैं: कुछ भी ठोस नहीं है।
\textbf{8.} \qty{1.0e-6}{mol/L}।
\textbf{9.} $14.00 - \frac13(38.55 - 6.00) = 3.15$।
\textbf{10.} $\log[\ce{Fe^3+}] = -38.55 + 3(14.00 - \mathrm{pH}) = 3.45 -
3\,\mathrm{pH}$: प्रवणता $-3$।
\textbf{11.} ज़िंक हाइड्रॉक्साइड के आरंभ तक, pH 6.96।
\textbf{12.} 3.15 से 6.96 तक।
\textbf{13.} \ce{Fe^3+ + OH- <=> FeOH^2+}, $\beta_1 =
[\ce{FeOH^2+}]/([\ce{Fe^3+}][\ce{OH-}]) = 10^{11.82}$।
\textbf{14.} दोनों बराबर होते हैं जब $[\ce{OH-}] = 10^{-11.82}$, अर्थात् pH 2.18 पर:
नीचे \ce{Fe^3+} प्रभावी है, ऊपर \ce{FeOH^2+}।
\textbf{15.} $10^{11.82 + 3.15 - 14.00} = 10^{0.97} = 9.3$: pH 3.15 पर
घुला आयरन मुक्त \ce{Fe^3+} का 10.3 गुना है, \qty{1.0e-5}{mol/L}:
केवल \qty{99.0}{\percent} हटता है। पहला अनुमान सुरक्षित नहीं था।
\textbf{16.} $[\ce{Fe}]_{\mathrm{dissolved}} = 10^{3.45 - 3\,\mathrm{pH}}\,
(1 + 10^{\mathrm{pH} - 2.18})$।
\textbf{17.} $\log[\ce{FeOH^2+}] = 11.82 + \log[\ce{Fe^3+}] + \mathrm{pH} -
14.00 = 1.27 - 2\,\mathrm{pH}$: प्रवणता $-2$।
\textbf{18.} $1.27 - 2\,\mathrm{pH} \approx -6.00$ से 3.64 मिलता है;
\ce{Fe^3+} पद को शामिल करने पर $10^{3.45 - 10.92}(1 + 10^{1.46}) =
\num{1.0e-6}$, $\mathrm{pH} = 3.64$ पर।
\textbf{19.} $[\ce{Zn^2+}] = \num{5.0e-5}$: $14.00 - \frac12(16.38 - 4.30) =
7.96$। मैग्नीशियम हाइड्रॉक्साइड 9.22 पर शुरू होता है: नहीं।
\textbf{20.} \ce{Zn(OH)2(s) + 2OH- <=> Zn(OH)4^2-}, $K = \beta_4 K_s =
10^{14.45 - 16.38} = 10^{-1.93}$।
\textbf{21.} $[\ce{OH-}]^2 = \num{5.0e-3}/10^{-1.93}$, $\mathrm{pH} =
13.81$। संचालक को क्षारक अधिक नहीं डालना चाहिए: pH 9.2 से ऊपर
मैग्नीशियम ज़िंक के साथ अवक्षेपित हो जाता है, और उससे बहुत ऊपर ज़िंक फिर से
विलयन में चला जाता है।
\textbf{22.} $\num{1.0e-3} \times 1000 \times 0.999 \times 106.8 =
\qty{107}{g}$ प्रति घन मीटर।
\textbf{23.} $\num{5.0e-3} \times 1000 \times 0.99 \times 99.4 =
\qty{492}{g}$ प्रति घन मीटर।
\textbf{24.} एक साथ 6.96 से ऊपर ले जाने पर दोनों हाइड्रॉक्साइड एक ही
गाद में अवक्षेपित होते हैं: ज़िंक आयरन के अपशिष्ट में खो जाता है, और अपशिष्ट
ज़िंक से दूषित हो जाता है। सीमाओं के बीच रुकने से केवल आयरन
हटता है।
\textbf{25.} \textbf{pH 3.64 से pH 6.96 तक}: 3.64 से नीचे आयरन का
\qty{0.1}{\percent} से अधिक अब भी घुला रहता है, अधिकतर \ce{FeOH^2+} के रूप में;
6.96 से ऊपर ज़िंक हाइड्रॉक्साइड बनता है। \ce{FeOH^2+} की उपेक्षा करने पर
निचली सीमा 3.15 पर आती।
\end{solution}
